% --------------------------
% CONFIGURATION
% --------------------------
\documentclass[11pt,a4paper]{moderncv}
\moderncvstyle{classic}
\moderncvcolor{black}
\definecolor{firstnamecolor}{rgb}{0,0,0}

% Basic packages
\usepackage[utf8]{inputenc}
\usepackage[T1]{fontenc}
\usepackage[scale=0.85]{geometry}
\usepackage[dvipsnames]{xcolor}
\definecolor{PaleBlue}{RGB}{76, 151, 215}
\usepackage[colorlinks=true, urlcolor=PaleBlue, linkcolor=PaleBlue]{hyperref}
\usepackage{microtype}

% Configure font settings
\renewcommand{\rmdefault}{ppl}
\renewcommand{\sfdefault}{phv}
\renewcommand*{\namefont}{\fontsize{26}{18}\mdseries}

% Configure microtype
\microtypesetup{expansion=false}

% Load sensitive information after copying into outputs/<listing_slug>/.
\input{../../secrets/personal_info_dk.tex}

% Title document
\title{Curriculum Vitae}

% %--------------------------

\begin{document}

\makecvtitle

\section*{Introduction}
Data scientist and ML engineer with an MSc in Engineering Physics and seven years of experience spanning research and production ML. My background combines deep learning, computer vision, and probabilistic modelling, including custom neural networks and representation learning from 3D medical images. I have led projects from problem formulation and experimentation through deployment, with an emphasis on understanding model behaviour and evaluating results critically. I am keen to apply this experience to challenging problems in environmental monitoring.


% ---------------------------
% Selected Work Experience
% ---------------------------
\section{Selected Experience}

% ---------------------------
\large{Researcher within Data Science @ RaySearch Laboratories}

\normalsize

\begin{itemize}
    \item Developed autoencoder-based patient embeddings and learned patient similarity metrics for end-to-end automated cancer treatment planning.
    \item Combined computer vision with Gaussian processes to predict probabilistic spatial dose distributions.
    \item Worked with medical physicists to define requirements and evaluate methods in safety-critical settings.

\end{itemize}

{\footnotesize\textit{Keywords: deep learning, 3D image analysis, representation learning, Gaussian processes, probabilistic modelling, optimisation}}

\hfill{\small{\textit{\hyperref[sec:raysearch]{Read more...}}}}

% ---------------------------
\large{Principal Data Scientist @ Valcon}
\normalsize
\begin{itemize}
    \item Led ML projects from problem formulation and experimentation through production deployment.
    \item Designed and deployed a custom, modular, real-time neural network using transfer learning with textual and high-dimensional categorical inputs.
    \item Mentored junior data scientists in ML development and software engineering.
\end{itemize}

{\footnotesize\textit{Keywords: deep learning, transfer learning, model architecture, production ML, technical leadership}}

\hfill{\small{\textit{\hyperref[sec:valcon]{Read more...}}}}


% ---------------------------

\large{\texttt{iris} -- Personal Computer Vision Project}
\normalsize
\begin{itemize}
    \item Developed a CV system to detect products in webshop images and match them to catalogue items.
    \item Fine-tuned YOLOv8 and combined OpenCLIP image embeddings with vector search for retrieval.
\end{itemize}

{\footnotesize\textit{Keywords: computer vision, object detection, transfer learning, representation learning, vector similarity search, ml engineering}}

\hfill{\small{\textit{\href{https://github.com/ivar-cashin-eriksson/iris}{GitHub}}}; \quad \small{\textit{\hyperref[sec:iris]{Read more...}}}}

% ---------------------------
% Education
% ---------------------------
\section{Education}

\cventry{2015--2020}{KTH, Royal Institute of Technology -- MSc. in Engineering Physics}{}{GPA: 5.0 / 5.0}{}{
    Specialisation: Applied Mathematics, Statistical Learning and Data Analytics. 
    \newline Relevant coursework: Statistical Machine Learning, Probability Theory, Optimisation. 
    \newline Thesis: \href{https://kth.diva-portal.org/smash/get/diva2:1431327/FULLTEXT02.pdf}{Probabilistic Dose Prediction for Radiation Therapy Planning}}
\cventry{2019}{ETH Zürich, Swiss Federal Institute of Technology -- Applied Mathematics}{Exchange}{}{}{}

\cventry{2017--2022}{SU, Stockholm University -- BSc. in Economics}{Additional Bachelor's Degree}{}{}{}

\cventry{2017}{ZJU, Zhejiang University}{Research exchange}{}{Laser trimming of microring resonators.}{}


% ---------------------------
% Skills and Other
% ---------------------------
\section{Skills and Other}

\normalsize
\cvitem{Technology}{Python, PyTorch, TensorFlow, scikit-learn, Git, Linux, Docker, Databricks, MLflow, SQL}
\cvitem{Methods}{Deep learning, computer vision, probabilistic modelling, representation learning}
\cvitem{Language}{Native English and Swedish, intermediate Danish}

\newpage

% ---------------------------
% Detailed Work Experience
% ---------------------------
\section{Detailed Work Experience}

% ---------------------------
\phantomsection\label{sec:signum}
\cventry{2024--Present}{Data Scientist / ML Engineer}{Signum Life Science}{Copenhagen}{}{
    At Signum, I work across machine learning and software engineering, translating loosely defined problems in pharmaceutical analytics into testable methods and maintainable production systems.
    \newline\hspace*{2em}
    I develop embedding-based retrieval and matching workflows for domain-specific text. This involves designing evaluation methods, analysing semantic similarity and retrieval failures, and iterating on modelling approaches based on observed errors.
    \newline\hspace*{2em}
    I also lead the architecture and development of a forecasting platform for pharmaceutical prices, demand, and patient populations. This involves defining modelling approaches, designing reproducible training and evaluation workflows, and building an architecture supporting thousands of independently trained models. The work combines statistical modelling, software engineering, and close collaboration with technical and domain stakeholders.
}{}
\vspace*{8pt}

% ---------------------------
\phantomsection\label{sec:valcon}
\cventry{2022--2024}{Principal Data Scientist}{Valcon}{Copenhagen}{}{
    At Valcon, I led ML projects across energy, IoT, pharmaceutical, and manufacturing domains, taking projects from problem formulation and model development through production deployment. I also mentored junior colleagues and delivered training.
    \newline\hspace*{2em}
    For one project, I designed and deployed a custom neural network for real-time prediction. Using transfer learning, the model combined natural-language and high-dimensional categorical inputs to predict high-dimensional categorical outputs. I was responsible for the model architecture, training and evaluation pipelines, and production integration.
    \newline\hspace*{2em}
    I received the Valcon People's Choice Award for exemplary work, leadership, team spirit, and scientific curiosity.
}{}
\vspace*{8pt}

% ---------------------------
\phantomsection\label{sec:valtech}
\cventry{2021--2022}{Data Scientist}{Valtech}{Copenhagen}{}{
    At Valtech, I developed ML solutions for e-commerce, including time-series forecasting and customer-behaviour modelling. I built forecasting pipelines using Prophet and applied clustering to behavioural data to support personalisation and customer segmentation.
    \newline\hspace*{2em}
    I worked with technical and commercial stakeholders to translate modelling results into operational recommendations and presented findings to senior decision-makers.
}{}
\vspace*{8pt}

% ---------------------------
\phantomsection\label{sec:raysearch}
\cventry{2019--2021}{Researcher within Data Science}{RaySearch Laboratories}{Stockholm}{}{
    At RaySearch Laboratories, I researched ML and optimisation methods for automated radiotherapy planning. The work combined deep learning, probabilistic modelling, and multi-objective optimisation in a safety-critical clinical setting.
    \newline\hspace*{2em}
    I developed a pipeline that learned latent anatomical representations from segmented 3D CT scans using autoencoders. These representations were used to identify similar patients through a learned similarity metric and to predict dose-volume histograms and spatial dose distributions using sparse Gaussian processes.
    \newline\hspace*{2em}
    I also redesigned lexicographic optimisation algorithms to better represent clinical decision hierarchies. I collaborated with medical physicists and clinicians to define requirements, evaluate model behaviour, and ensure that the resulting methods aligned with oncological standards and practical workflows.
}{}
\vspace*{8pt}

% ---------------------------
% Other Projects
% ---------------------------
\section{Personal Projects}
\cvitem{\phantomsection\label{sec:iris}Project \texttt{iris}}{
    Independently developed a computer vision system that detects products in webshop images and matches them to catalogue items. The system combines YOLOv8 object detection, OpenCLIP embeddings, and Qdrant vector similarity search.
    \newline\hspace*{2em}
    Built the complete application, including data collection and processing pipelines, a FastAPI backend, and a Chromium extension that displays product links directly within images. GitHub repository can be found \href{https://github.com/ivar-cashin-eriksson/iris}{here}.
}

\end{document}
